\documentclass[11pt,a4paper]{coverletter}

\senderblock
  {Radheem Bin Razi}
  {Distributed Systems \& Cloud-Native Engineer}
  {Ilmenau, Germany \textbar\ \href{mailto:sheikh.radheem@gmail.com}{sheikh.radheem@gmail.com} \textbar\ \href{https://radheem.github.io/my-cv}{portfolio}}

\begin{document}

%% ===== ENGLISH =====
\recipient{TU Ilmenau, Electronic Measurements and Signal Processing (EMS) Group / Prof. Thomas Dallmann}{Hiring Team}{}

\opening{Dear Hiring Team,}

\vspace{8pt}\noindent\textbf{1. Why TU Ilmenau?}\par\vspace{3pt}

Making a communication network perceive its environment turns the infrastructure into a sensor. Bridging that physical awareness with software-defined control is the engineering challenge I want to tackle. I am applying for the Master's Thesis Opportunity in the SENSATION project to apply my background in distributed systems and ML infrastructure to 6G integrated sensing and communication.

\vspace{8pt}\noindent\textbf{2. Why me?}\par\vspace{3pt}

\bullets

  \item Deployed an end-to-end O-RAN AIML framework on Kubernetes with Helm, automating the full ML lifecycle from feature groups to inference via KServe; achieved a German grade of 1.0 for this 15-credit research project.

  \item Built high-throughput Go microservices with gRPC and NATS JetStream to replace legacy orchestration, delivering the low-latency event-driven pipelines required for real-time cyber-physical coordination.

  \item Developed a config-driven Python SDK to automate ML model deployment and training, demonstrating the ability to build the software-defined instrumentation and automation layers needed for perceptive communication networks.

\bulletsend

I am available full-time immediately and also open to working-student or part-time arrangements; being based in Germany allows me to relocate to Ilmenau within two to three weeks.

\closing{Sincerely,}{Radheem Bin Razi}

\clearpage
\selectlanguage{ngerman}

%% ===== DEUTSCH =====
\recipient{TU Ilmenau, Electronic Measurements and Signal Processing (EMS) Group / Prof. Thomas Dallmann}{Personalabteilung}{}

\opening{Sehr geehrtes Hiring-Team,}

\vspace{8pt}\noindent\textbf{1. Warum TU Ilmenau?}\par\vspace{3pt}

Ein Kommunikationsnetzwerk seine Umgebung wahrnehmen zu lassen, verwandelt die Infrastruktur in einen Sensor. Die Verbindung dieser physischen Wahrnehmung mit einer softwaredefinierten Steuerung ist die ingenieurtechnische Herausforderung, die ich angehen möchte. Ich bewerbe mich um die Masterarbeit im Rahmen des SENSATION-Projekts, um mein Fachwissen in verteilten Systemen und ML-Infrastruktur auf die integrierte Sensorik und Kommunikation von 6G anzuwenden.

\vspace{8pt}\noindent\textbf{2. Warum ich?}\par\vspace{3pt}

\bullets

  \item Ein end-to-end O-RAN AIML-Framework auf Kubernetes mit Helm bereitgestellt, das den gesamten ML-Lebenszyklus von Feature-Gruppen bis zur Inferenz über KServe automatisiert; hierfür wurde in diesem 15-credit-Forschungsprojekt die deutsche Note 1,0 erzielt.

  \item Hochdurchsatz-Go-Mikrodienste mit gRPC und NATS JetStream entwickelt, um Legacy-Orchestrierung zu ersetzen und die für die Echtzeit-Koordination cyber-physischer Systeme erforderlichen, latenzarmen ereignisgesteuerten Pipelines bereitzustellen.

  \item Ein konfigurationsgetriebenes Python-SDK zur Automatisierung von ML-Modellbereitstellung und -Training entwickelt, das die Fähigkeit unterstreicht, die für wahrnehmungsfähige Kommunikationsnetze erforderlichen softwaredefinierten Instrumentierungs- und Automatisierungsschichten aufzubauen.

\bulletsend

Ich stehe ab sofort vollzeitlich zur Verfügung und bin auch an einer Werkstudententätigkeit oder Teilzeitbeschäftigung interessiert; da ich mich bereits in Deutschland befinde, kann ich innerhalb von zwei bis drei Wochen nach Ilmenau umziehen.

\closing{Mit freundlichen Grüßen,}{Radheem Bin Razi}

\end{document}
